\begin{table*}[t]
\centering
\caption{Association mapping of the focal free amino acids (FAAs) in the Goodman--Buckler panel: the two most abundant (asparagine, proline), the most heritable (lysine, asparagine, aspartate, proline), the amide glutamine and the most variable (methionine). Median, coefficient of variation (CV), line-mean broad-sense heritability ($H^2$) and genomic heritability ($h^2_g$) are from Supplementary Table~S3. Loci: loci (Supplementary Table~S7) significant at 5\% Benjamini--Hochberg FDR for at least one phenotype reporting on the amino acid (its absolute abundance and its ratios); in parentheses, loci for the absolute trait alone. Best $P$: smallest $P$ among those phenotypes. Trait-matched candidate genes: B73 RefGen\_v5 genes within 25~kb of a significant SNP for one of those phenotypes whose annotation places them in the metabolism of the amino acid or its biosynthetic family, in amino acid or nitrogen transport, or in primary nitrogen assimilation (Methods); for each, the phenotype with the smallest $P$, that $P$, and the distance from the SNP to the gene are given. $^{\dagger}$Significant in both BLINK and FarmCPU.}
\label{tab:amino_gwas_focus}
\footnotesize
\setlength{\tabcolsep}{3pt}
\begin{tabularx}{\textwidth}{@{}l r r r r r r >{\raggedright\arraybackslash}X@{}}
\hline
FAA & Median & CV (\%) & $H^2$ & $h^2_g$ & Loci & Best $P$ & Trait-matched candidate genes (annotation; phenotype, $P$, distance) \\
\hline
Asparagine (N) & 3.783 & 70 & 0.83 & 0.68 & 356 (73) & $2.8\times10^{-20}$ & \textit{Zm00001eb368680} (amino acid transporter AVT6C; \texttt{N}, $1.3\times10^{-8}$$^{\dagger}$, 10.3 kb); \textit{Zm00001eb329710} (glutamine amidotransferase type 2; \texttt{N.E}, $6.1\times10^{-10}$, 1.0 kb); \textit{Zm00001eb405410} (homoserine dehydrogenase; \texttt{N.IMNTDK}, $3.0\times10^{-9}$, in gene); \textit{Zm00001eb036060} (asparagine synthetase domain; \texttt{N.E}, $6.0\times10^{-7}$, 2.2 kb); \textit{Zm00001eb036070} (asparagine synthetase domain; \texttt{N.E}, $6.0\times10^{-7}$, 1.6 kb); \textit{Zm00001eb389290} (dihydrodipicolinate synthase; \texttt{N.Q}, $3.7\times10^{-7}$, 2.7 kb); \textit{Zm00001eb176170} (dihydrodipicolinate reductase; \texttt{N.E}, $6.1\times10^{-7}$, 4.3 kb) \\
Proline (P) & 3.202 & 63 & 0.80 & 0.40 & 59 (59) & $2.7\times10^{-15}$ & \textit{Zm00001eb217220} (glutamine amidotransferase; \texttt{P}, $2.2\times10^{-9}$, in gene) \\
Lysine (K) & 0.272 & 90 & 0.88 & 0.73 & 116 (54) & $8.3\times10^{-16}$ & \textit{Zm00001eb277160} (NRT1/PTR family 4.6; \texttt{K.Total}, $2.5\times10^{-11}$$^{\dagger}$, in gene); \textit{Zm00001eb360480} (glutamate synthase (NADH); \texttt{K.IMNTDK}, $2.8\times10^{-11}$, 15.1 kb); \textit{Zm00001eb352400} (WAT1-related (UmamiT); \texttt{K.IMNTDK}, $2.1\times10^{-8}$$^{\dagger}$, 3.6 kb) \\
Aspartate (D) & 1.178 & 76 & 0.81 & 0.84 & 91 (20) & $3.3\times10^{-16}$ & \textit{Zm00001eb207630} (WAT1-related (UmamiT); \texttt{D.IMNTDK}, $2.7\times10^{-10}$, 0.4 kb) \\
Glutamine (Q) & 0.371 & 130 & 0.69 & 0.85 & 165 (65) & $4.9\times10^{-22}$ & \textit{Zm00001eb379220} (WAT1-related (UmamiT); \texttt{Q}, $1.2\times10^{-11}$, in gene); \textit{Zm00001eb261470} (amino acid transporter; \texttt{Q.E}, $8.4\times10^{-8}$, in gene); \textit{Zm00001eb261480} (amino acid transporter; \texttt{Q.E}, $8.4\times10^{-8}$, 6.5 kb); \textit{Zm00001eb261490} (amino acid transporter; \texttt{Q.E}, $8.4\times10^{-8}$, 13.5 kb) \\
Methionine (M) & 0.047 & 180 & 0.78 & 0.46 & 104 (56) & $3.2\times10^{-17}$ & none \\
\hline
\end{tabularx}
\end{table*}
